\subsection{展开存储与矩阵分块复用}
\label{l11:subsec:unroll_cost}

显式展开把窗口重叠转化为数组中的重复元素。全部输入共有 $BCHW$ 个元素，全部展开矩阵共有 $BDL=BCK^2H_{\mathrm{o}}W_{\mathrm{o}}$ 个元素，因此展开数组相对于输入数组的大小比为
\begin{equation}
  \rho_{\mathrm{expand}}
  =\frac{BCK^2H_{\mathrm{o}}W_{\mathrm{o}}}{BCHW}
  =\frac{K^2(H-K+1)(W-K+1)}{HW}.
  \label{l11:eq:expansion_ratio}
\end{equation}
式\eqref{l11:eq:small_unroll}中的展开矩阵有 $12\times4=48$ 个元素，原输入有 $3\times3\times3=27$ 个元素，大小比为 $48/27=16/9\approx1.78$。当 $H,W$ 远大于 $K$ 时，$H_{\mathrm{o}}W_{\mathrm{o}}$ 接近 $HW$，该比值趋近 $K^2$。这不是每个输入元素都复制相同次数：内部元素可能属于许多窗口，边界元素属于较少窗口；式\eqref{l11:eq:expansion_ratio}比较的是整个数组的大小。

若每个元素占 $s$ 字节，完整展开工作区需要 $sBDL$ 字节。代码\ref{l11:lst:unroll_cpu}或\ref{l11:lst:unroll_gpu}对每个展开位置读取一次原输入并写入一次展开数组，因而产生 $BDL$ 次输入读取请求和 $BDL$ 次展开写入；随后矩阵乘法还要读取展开数组。缓存、事务合并和矩阵分块会影响实际显存传输量，因此这里的元素请求计数不能直接等同于硬件上的总传输字节数。

展开并没有减少卷积的乘加次数。对每幅图像计算一个 $M\times D$ 与 $D\times L$ 的矩阵乘积，仍需 $MDL$ 次乘加，整个批次为 $BMDL$ 次。其优势在于把输入、权重和输出之间的复用关系组织成适合通用矩阵乘法（general matrix multiplication，GEMM）的规则形式。

考虑矩阵乘法中的一个输出块，它覆盖 $T_m$ 个输出通道和 $T_n$ 个展开列；一次沿求和维度处理 $T_k$ 个元素。这里两个输出方向是通道 $m$ 和线性化空间位置 $\ell$，并非直接卷积空间块中的高度和宽度。在这一阶段，线程块将 $T_m\times T_k$ 的核矩阵块与 $T_k\times T_n$ 的展开矩阵块放入共享内存，再计算它们对输出块的贡献。对于三个方向均完整有效的分块，加载元素数与乘加数分别为
\begin{equation}
  N_{\mathrm{load}}=T_mT_k+T_kT_n,
  \qquad
  N_{\mathrm{MAC}}=T_mT_nT_k.
  \label{l11:eq:gemm_tile_reuse}
\end{equation}
每个核元素可用于 $T_n$ 个输出位置，每个展开元素可用于 $T_m$ 个输出通道。相比只让一个线程块负责单一输出通道的空间块，矩阵分块还能够在块内显式复用不同输出通道需要的同一输入窗口。完成全部求和块后，累加结果才写回对应的输出区域。

输出块更大通常带来更多块内复用，也需要更多共享内存、累加寄存器和线程工作量。边缘块中不足 $T_m,T_n,T_k$ 的部分还会产生无效位置。选择块尺寸时，需要同时考虑有效工作比例和可驻留的线程块数量，不能只追求单个块的复用次数。

矩阵乘法的效率还取决于三个维度 $M,D,L$ 各自的大小。一个维度很大，并不能补偿另一个维度中的大量无效分块位置。例如，第一层卷积的矩阵形状为 $4\times49$ 与 $49\times6400$：空间位置很多，但输出通道只有 $4$ 个。若沿输出通道方向固定使用 $T_m=16$ 的分块，则该方向只有四分之一的位置有效；若求和分块宽度为 $T_k=16$，最后一个求和块也只有一个有效元素。调整分块形状或让线程承担不同数量的输出，可以减少这些损失，但会改变资源占用和访存组织。较小的矩阵还可能使内核启动开销占据较大比例。因此，将卷积写成 GEMM 提供了规则的复用结构，实际性能仍需结合矩阵形状、边缘利用率和完整执行阶段衡量。

\subsection{隐式展开与共享内存、寄存器联合分块}
\label{l11:subsec:implicit_unroll}

隐式展开（implicit unrolling）保留 $V_b=AU_b$ 的矩阵计算方式，但不在全局内存中存放完整 $U_b$。当分块矩阵乘法准备加载展开矩阵中的某个元素时，直接利用式\eqref{l11:eq:unroll_inverse}算出它在原输入中的地址，将原输入值送入当前共享内存块。展开矩阵只作为索引关系存在。

固定图像 $b$，设当前展开矩阵块的左上位置为 $(r_0,\ell_0)$，块内位置为 $(k,j)$。首先得到 $r=r_0+k$、$\ell=\ell_0+j$；若 $r<D$ 且 $\ell<L$，就从 $(r,\ell)$ 恢复 $c,p,q,h,w$，再读取 $X_{b,c,h+p,w+q}$。核矩阵块中的元素仍是 $A_{m_0+i,r_0+k}$。两块就绪后，输出块的每个元素累加
\begin{equation}
  \operatorname{acc}_{i,j}
  \mathrel{+}=\sum_{k=0}^{T_k-1}
  A_{m_0+i,r_0+k}\,(U_b)_{r_0+k,\ell_0+j}.
  \label{l11:eq:implicit_tile}
\end{equation}
超出 $M,D,L$ 范围的分块位置填零，最终只写回合法输出。这里的零用于补齐矩阵边缘块，并未给原输入增加卷积填充；合法输出仍然使用原输入内部的完整窗口。共享内存块的加载、使用和下一轮覆盖之间，同样需要正确的块内同步。

把展开地址计算并入矩阵乘法，是核函数融合（kernel fusion）的一种形式。它省去了完整展开数组的全局写入和独立展开阶段的启动开销；原先读取 $U$ 的动作改为按需读取 $X$，无需再从全局内存读取物化的展开数组，同时也省去了 $sBDL$ 字节的完整展开工作区。共享内存中的局部块与输出累加寄存器仍然需要空间。原输入的一个元素可能因属于多个窗口或多个线程块而再次被读取，隐式展开不会自动消除所有重叠读取。它还增加了加载路径中的取整、取余和地址计算，因此是否更快取决于减少的存储开销与新增的寻址成本。

在共享内存分块之上，还可以让每个线程保存一个 $R_m\times R_n$ 的输出子块。对一个求和位置，该线程从共享内存取出 $R_m$ 个权重和 $R_n$ 个输入值，组成两组短向量；它们的所有两两乘积更新 $R_mR_n$ 个寄存器累加器。这样，一次取到的权重可在本线程内用于 $R_n$ 个输出位置，一次取到的输入可用于 $R_m$ 个输出通道。线程块之间的分工、共享内存块和线程内寄存器子块由此形成两级复用。

寄存器子块增大后，每个线程要保存更多累加器和临时值，可能减少同时驻留的线程，过高的寄存器需求还可能造成溢出存储。因此，共享内存与寄存器联合分块应平衡数据复用和执行资源占用。四种实现的主要区别可归纳为表\ref{l11:tab:implementation_routes}；联合分块既可用于显式展开后的矩阵乘法，也可用于隐式展开。

\begin{table}[htbp]
\centering
\small
\caption{卷积实现方式与主要数据组织}
\label{l11:tab:implementation_routes}
\begin{tabularx}{\linewidth}{L{2.35cm}XX}
\toprule
实现方式 & 数据复用方式 & 主要代价或约束\\
\midrule
直接空间分块 & 协作装入带晕区的输入；适合时广播权重 & 晕区及不同块的重复读取；通道间同步\\
显式展开后乘法 & 完整保存窗口列，再用矩阵分块共享输入和权重 & 展开工作区、展开读写与独立执行阶段\\
隐式展开乘法 & 装入矩阵块时从原输入取值 & 地址计算；输入访问仍可能重复\\
共享内存与寄存器联合分块 & 块内共享数据，线程内复用短向量 & 累加寄存器与共享内存占用\\
\bottomrule
\end{tabularx}
\end{table}

这些实现计算的是同一组实数乘积与求和，但浮点加法不满足任意重排下的结合律。分块、并行归约以及融合乘加指令可能改变舍入位置或累加顺序，因此数值核验应同时检查布局、边界和合理的浮点误差，不能把所有非零差异都判定为索引错误。
